\documentclass[aspectratio=169,xcolor=dvipsnames]{beamer}

\AtBeginSection[]
{
    \begin{frame}
        \frametitle{Visão Geral}
        \tableofcontents[currentsection]
    \end{frame}
}

\usefonttheme{professionalfonts} % using non standard fonts for beamer
\usefonttheme{serif} % default family is serif
\usepackage{fontspec}
\setmainfont{XCharter}[
    UprightFont    = *-Roman,
    BoldFont       = *-Bold,
    ItalicFont     = *-Italic,
    BoldItalicFont = *-BoldItalic,
]

\usepackage{hyperref}
\usepackage{graphicx} % Allows including images
\usepackage{booktabs} % Allows the use of \toprule, \midrule and \bottomrule in tables
\input{./slides/SimplePlusTheme/beamerthemeSimplePlus.sty}
\input{./slides/SimplePlusTheme/beamerinnerthemeSimplePlus.sty}
\input{./slides/SimplePlusTheme/beamerfontthemeSimplePlus.sty}
\input{./slides/SimplePlusTheme/beamercolorthemeSimplePlus.sty}

\usepackage{listings}
\setbeamercovered{transparent}

\usepackage{biblatex}
\addbibresource{slides/reference.bib}
\renewcommand*{\bibfont}{\footnotesize}
% Tamanho dos slides de "Leituras Recomendadas": maior que a bibliografia
% completa (\bibfont acima, que rende \tiny via a redefinição de \footnotesize)
% para destacar as leituras curadas.
\newcommand{\leiturasize}{\small}

\usepackage[brazilian ]{babel}
\usepackage{csquotes}
\usepackage{pgfplots}
\pgfplotsset{compat=1.18}

\usepackage{emoji}
\usepackage{amsmath}
\usepackage[xcharter,bigdelims,vvarbb]{newtxmath}
% newtxmath's operators font requests the fontspec family `XCharter(0)' in OT1;
% without these substitutions math digits and operator names silently fall back
% to Computer Modern (LaTeX Font Warning: Font shape `OT1/XCharter(0)/m/n' undefined).
\DeclareFontFamily{OT1}{XCharter(0)}{}
\DeclareFontShape{OT1}{XCharter(0)}{m}{n}{<->ssub * XCharter-TLF/m/n}{}
\DeclareFontShape{OT1}{XCharter(0)}{m}{it}{<->ssub * XCharter-TLF/m/it}{}
\DeclareFontShape{OT1}{XCharter(0)}{b}{n}{<->ssub * XCharter-TLF/b/n}{}
\DeclareFontShape{OT1}{XCharter(0)}{bx}{n}{<->ssub * XCharter-TLF/b/n}{}

\usepackage{bm}
\usepackage[table]{xcolor}
\usepackage{xcolor}
\usepackage{algpseudocode}

\usepackage{cancel}
\usepackage{ulem}

\usetikzlibrary{shapes}
\usetikzlibrary{arrows}
\usetikzlibrary {arrows.meta}
\usetikzlibrary{calc,positioning}
\usepgfplotslibrary{fillbetween}
\usetikzlibrary{angles,quotes}
\usetikzlibrary{tikzmark}


\definecolor{PastelRed}{HTML}{FFC1CC}
\definecolor{PastelBlue}{HTML}{C2F0FF}
\definecolor{PastelBlue2}{HTML}{3182BD}

\definecolor{PastelPink}{HTML}{FFCBE1}
\definecolor{PastelGreen}{HTML}{D6E5BD}
\definecolor{PastelYellow}{HTML}{F9E1A8}
\definecolor{PastelBlue}{HTML}{BCD8EC}
\definecolor{PastelPurple}{HTML}{DCCCEC}
\definecolor{PastelOrange}{HTML}{FFDAB4}


\renewcommand{\footnotesize}{\tiny}

\newcommand{\mespor}{
  \ifcase\month
    \or janeiro
    \or fevereiro
    \or março
    \or abril
    \or maio
    \or junho
    \or julho
    \or agosto
    \or setembro
    \or outubro
    \or novembro
    \or dezembro
  \fi
}
% \usepackage{csquotes}
% \usepackage{minted}
% \usemintedstyle{xcode}
\definecolor{vscLightBackground}{HTML}{FFFFFF}
\definecolor{vscLightText}{HTML}{000000}
\definecolor{vscLightGreen}{HTML}{008000}        % For Comments AND Numbers
\definecolor{vscLightRed}{HTML}{A31515}          % For Strings
\definecolor{vscLightTeal}{HTML}{267F99}
\definecolor{vscBlue}{HTML}{0000FF}% For main keywords (if, for, def...)
\definecolor{vscLightPytorch}{HTML}{800080}      % For PyTorch keywords (dark magenta)

\lstdefinestyle{vscode-light}{
    language=Python,
    backgroundcolor=\color{vscLightBackground},
    basicstyle=\ttfamily\scriptsize\color{vscLightText},
    keywordstyle=\color{vscLightTeal},
    commentstyle=\color{vscLightGreen},
    stringstyle=\color{vscLightRed},
    % Style for PyTorch keywords (using keyword group 2)
    keywordstyle=[2]\color{vscLightPytorch}, 
    morekeywords={self, True, False, None},
    morekeywords=[2]{
        torch, Tensor, device, dtype, to, nn, optim, autograd, utils, data,
        Module, Linear, Conv2d, ReLU, CrossEntropyLoss, Adam, SGD,
        Dataset, DataLoader, randn, zeros, ones, arange, cat, stack,
        view, reshape, sum, mean, max, min, matmul, no_grad, backward
    },
    keywordstyle=[3]\color{vscBlue}, 
    morekeywords=[3]{
        model, loss, y, y_hat, x , optimizer
    },
    frame=single,
    tabsize=1,
    showstringspaces=false,
    breaklines=true,
    captionpos=b,
    extendedchars=true,
    numbers=left,
}

\usepackage[cache=true]{minted}
% minted v3+ (TeX Live 2024+) auto-detects the output directory.
% minted v2 (Ubuntu apt TeX Live 2023) needs it set explicitly to
% match latexmkrc's $out_dir = 'build'.
\makeatletter
\@ifpackagelater{minted}{2024/01/01}{}{%
  \def\minted@outputdir{build/}%
}
\makeatother
\setminted{
    fontsize=\scriptsize,
    style=vs,
    breaklines=true,
    numbers=left,
    numbersep=5pt,
    frame=single,
}
\providecommand{\citeauthoryear}[1]{(\citeauthor{#1},~\citeyear{#1})}

\providecommand{\thetav}{\bm{\theta}}

% vetor coluna 2D com colchetes
\providecommand{\cvec}[2]{\begin{bmatrix} #1 \\ #2 \end{bmatrix}}

\providecommand{\varvector}[1]{\mathbf{#1}}
% instance
\providecommand{\X}{\varvector{X}}
\providecommand{\mlinstance}{\varvector{X}}
\providecommand{\mlinstancei}{\varvector{X}^{(i)}}
\providecommand{\mlinstancex}[1]{\varvector{X}^{(#1)}}

% target
\providecommand{\y}{\varvector{y}}
\providecommand{\mltarget}{\varvector{y}}
\providecommand{\mltargeti}{\varvector{y}^{(i)}}
\providecommand{\mltargetx}[1]{\varvector{y}^{(#1)}}

%hats
\providecommand{\yhat}{\hat{\varvector{y}}}
\providecommand{\mltargethat}{\hat{\varvector{y}}}
\providecommand{\mltargethati}{\hat{\varvector{y}}^{(i)}}
\providecommand{\mltargethatx}[1]{\hat{\varvector{y}}^{(#1)}}

% linear reg
\providecommand{\mllinearregression}{\mltargethat=\mlinstance\bm{\theta}}

\providecommand{\tikzarrowcolor}[3]{%
  \begin{tikzpicture}[remember picture, overlay]
    % Arrow body
    \path[fill=#3] ([shift={(#1,#2)}]current page.south west) 
      rectangle ++(0.3, -0.2);
    % Arrow head
    \path[fill=#3] ([shift={(#1+0.3,#2+0.1)}]current page.south west) -- 
         ++(0, -0.4) -- 
         ++(0.2, 0.2) -- 
         cycle;
  \end{tikzpicture}%
}

\providecommand{\tikzarrow}[2]{%
    \tikzarrowcolor{#1}{#2}{Red}
}



\DeclareMathOperator*{\argmax}{argmax}
\DeclareMathOperator*{\argmin}{argmin}


%------------------------------------------------------------------------
%   SHARED COLOR PALETTE FOR UNIT 05: GENERATIVE MODELS
%
%   Input this file (after preamble/custom-commands) in every lecture of
%   this unit so that the same concept always gets the same color.
%------------------------------------------------------------------------

% Encoder / Discriminator / "Going into the latent" path
\definecolor{EncRed}{HTML}{C62828}

% Decoder / Generator / "Coming out of the latent" path
\definecolor{DecBlue}{HTML}{1565C0}

% Aliases so either naming convention works (GANs call them Gen/Disc,
% VAEs/diffusion call them Encoder/Decoder).
\colorlet{GenRed}{EncRed}
\colorlet{DiscBlue}{DecBlue}

% Reconstruction / "true-data" anchor
\colorlet{ReconGreen}{MediumGreen}

% Noise / stochastic perturbation
\definecolor{NoiseAmber}{HTML}{F9A825}

% Accent for cross-references (prior lecture, forward pointers)
\definecolor{AccentBlue}{HTML}{1565C0}
\colorlet{FlowRed}{EncRed}
\colorlet{FlowGreen}{MediumGreen}

% Emphasis box color for pedagogical asides
\definecolor{SoftAmber}{HTML}{F9A825}

%------------------------------------------------------------------------
%   SHARED TIKZ: GENERATIVE-MODELING TRILEMMA
%
%   TikZ replacement for slides/img/05-geradores/generative-trilemma.png.
%   Input this file AFTER colors.tex in every lecture of unit 05.
%
%   Usage:  \resizebox{\linewidth}{!}{\generativetrilemma[<family>]}
%   where <family> in {gan, diffusion, vaenf, none} highlights the family
%   relevant to the current deck (default: none = neutral figure).
%------------------------------------------------------------------------

\usepackage{etoolbox} % \ifstrequal (no-op if beamer already loaded it)

\newcommand{\generativetrilemma}[1][none]{%
  % --- per-family stroke width and opacity, driven by the [#1] highlight ---
  \def\wGan{1pt}\def\oGan{1}%
  \def\wDif{1pt}\def\oDif{1}%
  \def\wVae{1pt}\def\oVae{1}%
  \ifstrequal{#1}{none}{}{%
    \def\oGan{0.30}\def\oDif{0.30}\def\oVae{0.30}%
    \ifstrequal{#1}{gan}{\def\oGan{1}\def\wGan{1.7pt}}{}%
    \ifstrequal{#1}{diffusion}{\def\oDif{1}\def\wDif{1.7pt}}{}%
    \ifstrequal{#1}{vaenf}{\def\oVae{1}\def\wVae{1.7pt}}{}%
  }%
  \begin{tikzpicture}[
      corner/.style={circle, draw=black, thick, minimum size=2.4cm,
                     align=center, font=\footnotesize},
      contour/.style={dashed, line cap=round},
    ]
    % --- three corner circles ---
    \node[corner] (T) at (0,2.4)   {High Quality\\Samples};
    \node[corner] (L) at (-2.5,-1.6) {Fast\\Sampling};
    \node[corner] (R) at (2.5,-1.6)  {Mode\\Coverage /\\Diversity};

    % --- family contours, each enclosing the two corners it reaches ---
    % GANs: quality + fast sampling (top + bottom-left)
    \draw[contour, DiscBlue, line width=\wGan, opacity=\oGan,
          rotate around={58:(-1.25,0.4)}]
          (-1.25,0.4) ellipse [x radius=3.9, y radius=1.5];
    % Diffusion: quality + coverage (top + bottom-right)
    \draw[contour, EncRed, line width=\wDif, opacity=\oDif,
          rotate around={-58:(1.25,0.4)}]
          (1.25,0.4) ellipse [x radius=3.9, y radius=1.5];
    % VAEs / Normalizing Flows: fast sampling + coverage (bottom edge)
    \draw[contour, NoiseAmber, line width=\wVae, opacity=\oVae]
          (0,-1.6) ellipse [x radius=4.0, y radius=1.5];

    % --- family labels, color-matched ---
    \node[DiscBlue, align=center, font=\small, opacity=\oGan]
      at (-4.9,2.2) {Generative\\Adversarial\\Networks};
    \node[EncRed, align=center, font=\small, opacity=\oDif]
      at (4.9,2.2) {Denoising\\Diffusion\\Models};
    \node[NoiseAmber, align=center, font=\small, opacity=\oVae]
      at (0,-3.6) {Variational Autoencoders,\\Normalizing Flows};
  \end{tikzpicture}%
}


%------------------------------------------------------------------------
%   LECTURE-LOCAL EXTRAS (only what the shared palette doesn't cover)
%------------------------------------------------------------------------
\definecolor{SoftGray}{HTML}{F3F4F6}

%------------------------------------------------------------------------
%   TITLE PAGE
%------------------------------------------------------------------------

\title{Aprendizado Profundo}
\subtitle{\textit{Generative Adversarial Networks}}

\author{Lucas Silveira Kupssinskü}

\institute
{
    Machine Learning Theory and Applications Laboratory \\
    PUCRS
}
\date{\mespor de \the\year}

\begin{document}

  \begin{frame}
    \titlepage
  \end{frame}

  \begin{frame}{Visão Geral}
    \tableofcontents
  \end{frame}

  %======================================================================
  % SECTION 1: O QUE SÃO GANS
  %======================================================================
  \section{O que são GANs}

  %--- Slide: Definição ---
  \begin{frame}{\textit{Generative Adversarial Network}\footfullcite{goodfellow2014generative}}
    Uma \textbf{GAN} é um modelo não supervisionado cujo objetivo é
    \textbf{gerar novas instâncias indistinguíveis} do conjunto de dados
    de treinamento.
    \vspace{0.5em}
    \begin{itemize}
      \item Não constrói explicitamente uma distribuição de
      probabilidade sobre os dados.
      \item Em vez disso, aprende a \emph{amostrar} dessa distribuição.
      \item Usualmente é formada por \textbf{duas redes} treinadas em
      conjunto:
      \begin{itemize}
        \item \textcolor{GenRed}{\textbf{Gerador}}: mapeia uma distribuição
        aleatória para o espaço dos dados.
        \item \textcolor{DiscBlue}{\textbf{Discriminador}} (ou crítico):
        classifica as instâncias como reais ou sintéticas.
      \end{itemize}
      \item Desenvolvida originalmente por Ian Goodfellow em 2014.
    \end{itemize}
    \note{
      Pontos importantes para a aula:
      \begin{itemize}
        \item GANs são modelos \emph{implícitos} de distribuição: não
        calculamos $P(x)$, apenas sabemos amostrar dela.
        \item A ideia central é transformar um problema de modelagem
        densa em um problema de classificação, algo que redes neurais
        resolvem bem.
        \item O termo \emph{adversarial} vem do fato de gerador e
        discriminador terem objetivos opostos e se treinarem em
        conjunto.
      \end{itemize}
    }
  \end{frame}

  %--- Slide: GAN pipeline visual (anchor para Gerador + Discriminador) ---
  \begin{frame}{Gerador e Discriminador em conjunto}
    \centering
    \includegraphics[width=0.9\linewidth]{slides/img/05-geradores/gan-pipeline.png}\\
    \vspace{0.3em}
    \tiny Adaptado de Prince (2023), \emph{Understanding Deep Learning},
    cap.~15.
    \note{
      Visão geral apresentada cedo para ancorar visualmente as duas
      redes. As fórmulas de loss aparecem no diagrama como
      \emph{teaser}, serão desenvolvidas a partir da seção
      \emph{Loss do Discriminador}.
    }
  \end{frame}

  %--- Slide: Aplicações ---
  \begin{frame}{Aplicações}
    Onde GANs já mostraram resultados marcantes:
    \begin{itemize}
      \item \textbf{Síntese de faces fotorrealistas}: \textit{StyleGAN}
      \footfullcite{karras2019style} produz retratos indistinguíveis de
      pessoas reais.
      \item \textbf{Image-to-image translation}:
      \textit{pix2pix}\footfullcite{isola2017image} (mapa $\to$ foto,
      esboço $\to$ foto) e \textit{CycleGAN}\footfullcite{zhu2017unpaired}
      (cavalo $\leftrightarrow$ zebra) sem pares alinhados.
      \item \textbf{Super resolução}: ampliar imagens preservando
      textura plausível.
      \item \textbf{Aumento de dados}: gerar amostras sintéticas para
      classes raras em domínios médicos e industriais.
      \item \textbf{Deepfakes}: vídeo e áudio sintéticos de pessoas
      reais, com todas as questões éticas que isso traz.
    \end{itemize}
    \note{
      Bom momento para uma rápida discussão sobre uso responsável.
      Mencionar que muitas dessas aplicações já migraram para modelos
      difusores nos últimos anos, mas o ferramental conceitual de GAN
      ainda é central.
    }
  \end{frame}

  %--- Slide: Princípios - processo gerativo ---
  \begin{frame}{\textit{GANs} -- Princípios}
    \textbf{Objetivo:} gerar novas instâncias
    $\{x^{(j)*}\}$ que sigam a mesma distribuição dos dados de treino
    $\{x^{(i)}\}$.
    \vspace{0.5em}
    \begin{block}{Processo de geração de uma nova instância}
      \begin{enumerate}
        \item Amostrar uma \textbf{variável latente}
        $z^{(j)}$ de uma distribuição simples (normal, uniforme, \ldots).
        \item Usar $z^{(j)}$ como entrada para uma rede neural:
        \[
          x^{(j)*} = g\!\left[z^{(j)}, \thetav\right]
        \]
        onde $\thetav$ são os parâmetros do gerador aprendidos durante o
        treino.
      \end{enumerate}
    \end{block}
    \note{
      O latente $z$ é tipicamente de baixa dimensionalidade (p.~ex.\ 100
      dimensões) enquanto $x$ pode ser uma imagem de milhares de pixels.
      O gerador realiza um \emph{upsampling} semântico.
    }
  \end{frame}

  %--- Slide: Princípios - indistinguibilidade ---
  \begin{frame}{\textit{GANs} -- Princípios}
    \begin{itemize}
      \item Teremos atingido o objetivo quando as instâncias geradas
      forem \textbf{indistinguíveis} das instâncias do conjunto de
      dados.
      \item Introduzimos então uma segunda rede
      $f\!\left[\cdot,\phi\right]$ que tenta \textbf{classificar} suas
      entradas como reais ou geradas.
      \begin{itemize}
        \item Se o discriminador consegue classificar bem, ele pode
        gerar um \emph{sinal de gradiente} usado para melhorar o
        gerador.
        \item Se o discriminador não consegue distinguir, o gerador
        atingiu o objetivo.
      \end{itemize}
    \end{itemize}
    \begin{alertblock}{Intuição de jogo}
      O gerador tenta enganar o discriminador; o discriminador tenta não
      ser enganado. Um equilíbrio de Nash nesse jogo corresponde a
      $P_{\text{gerada}} = P_{\text{real}}$.
    \end{alertblock}
    \vspace{0.3em}
    Para enxergar essa dinâmica, vamos olhar um exemplo unidimensional
    em que o gerador se reduz a um \emph{único parâmetro}.
  \end{frame}

  %--- Slide: Exemplo toy - parte 1 ---
  \begin{frame}{\textit{GANs} -- Princípios: um exemplo \textit{toy}}
    \begin{columns}[T]
      \begin{column}{0.48\textwidth}
        \begin{itemize}
          \item Os dados reais foram obtidos de uma normal
          $\mathcal{N}(7,1)$.
          \item O latente é amostrado de $z^{(j)} \sim \mathcal{N}(0,1)$
          e o modelo gerador possui \textbf{um único parâmetro} que
          translada essa distribuição:
          \[
            x^{(j)*} = z^{(j)} + \theta.
          \]
          \item A sigmóide representa o classificador binário ideal.
        \end{itemize}
      \end{column}
      \begin{column}{0.50\textwidth}
        \centering
        \includegraphics[width=\linewidth]{slides/img/05-geradores/sigmoid-panels.png}\\
        \tiny Painel (a) com $\theta=3.0$.
      \end{column}
    \end{columns}
    \note{
      Esse exemplo \emph{toy} tem uma única variável de decisão no
      gerador: a translação $\theta$. Ele é útil para visualizar por que
      o discriminador empurra o gerador para a região correta: onde as
      distribuições se sobrepõem, a sigmóide fica próxima de 0.5, e o
      sinal para o gerador desaparece.
    }
  \end{frame}

  %--- Slide: Exemplo toy - parte 2 ---
  \begin{frame}{\textit{GANs} -- Princípios: sobreposição de distribuições}
    \begin{itemize}
      \item Valores diferentes de $\theta$ fazem com que o classificador
      tenha maior ou menor dificuldade em separar as distribuições.
      \item Quando $\theta \approx 7$, as distribuições coincidem e o
      classificador ideal degenera para chance aleatória (sigmóide
      quase plana).
    \end{itemize}
    \centering
    \includegraphics[width=0.9\linewidth]{slides/img/05-geradores/sigmoid-panels.png}\\
    \tiny (a) distribuições separadas; (b) parcialmente sobrepostas;
    (c) quase coincidentes.
  \end{frame}

  %======================================================================
  % SECTION 2: MINMAX LOSS
  %======================================================================
  \section{MinMax Loss}

  %--- Slide: Discriminator loss ---
  \begin{frame}{Loss do Discriminador}
    Para o discriminador, temos um problema de \textbf{classificação
    binária}. Rotulamos as instâncias reais com $y=1$ e as geradas com
    $y=0$, e minimizamos a entropia cruzada:
    \vspace{0.4em}
    \[
      \hat{\phi} = \argmin_{\phi} \Bigg[ \sum_i
        -\!\left(1 - y^{(i)}\right)\ln\!\left[1 - \sigma\!\left(f\!\left[x^{(i)},\phi\right]\right)\right]
        - y^{(i)} \ln\!\left[\sigma\!\left(f\!\left[x^{(i)},\phi\right]\right)\right] \Bigg]
    \]
    \vspace{0.2em}
    \begin{itemize}
      \item $\sigma(\cdot)$ é a sigmóide que mapeia o logit $f[\cdot,\phi]$
      para probabilidade de ``real''.
      \item A primeira parcela penaliza rotular geradas como reais; a
      segunda penaliza rotular reais como geradas.
    \end{itemize}
  \end{frame}

  %--- Slide: Discriminator loss split real vs gen ---
  \begin{frame}{Loss do Discriminador (separando reais e geradas)}
    Para tornar mais evidente que algumas instâncias são geradas
    ($y^{(j)}=0$), reescrevemos a loss em duas somas distintas:
    \vspace{0.3em}
    \[
      \hat{\phi} = \argmin_{\phi} \Bigg[ \sum_{j}
        \!\!\left(-\ln\!\left[1 - \sigma\!\left(f\!\left[\textcolor{GenRed}{x^{(j)*}},\phi\right]\right)\right]\right)
        - \sum_{i}
        \!\!\left(\ln\!\left[\sigma\!\left(f\!\left[\textcolor{DiscBlue}{x^{(i)}},\phi\right]\right)\right]\right) \Bigg]
    \]
    \vspace{0.2em}
    \begin{itemize}
      \item $j$ indexa instâncias \textcolor{GenRed}{geradas}
      ($x^{(j)*}$), $i$ indexa instâncias
      \textcolor{DiscBlue}{reais} ($x^{(i)}$).
      \item O gerador entra apenas no primeiro somatório, por meio das
      amostras $x^{(j)*} = g[z^{(j)},\thetav]$.
    \end{itemize}
  \end{frame}

  %--- Slide: MinMax formulation ---
  \begin{frame}{\textit{MinMax Loss}}
    Plugando a definição do gerador e \textbf{invertendo o objetivo do
    gerador}, obtemos a formulação adversarial conjunta:
    \vspace{0.3em}
    \[
      \hat{\phi},\hat{\thetav} =
      \argmin_{\phi}\, \argmax_{\thetav}
      \Bigg[ \sum_{j}\!\!
        \left(-\ln\!\left[1-\sigma\!\left(f\!\left[g[z^{(j)},\thetav],\phi\right]\right)\right]\right)
        - \sum_{i}\!\!
        \left(\ln\!\left[\sigma\!\left(f\!\left[x^{(i)},\phi\right]\right)\right]\right) \Bigg]
    \]
    \vspace{0.3em}
    \begin{block}{\textit{MinMax Loss}}
      O discriminador \textbf{minimiza} a entropia cruzada (classifica
      bem) enquanto o gerador \textbf{maximiza} o mesmo objetivo
      (tenta enganar o discriminador).
    \end{block}
    \note{
      Observação prática: ao invés de ``maximizar'' a loss do
      discriminador, o gerador normalmente minimiza
      $-\ln \sigma(f[g(z),\phi])$ (\emph{non-saturating loss}),
      que tem gradiente melhor quando o discriminador está vencendo o
      jogo.
    }
  \end{frame}

  %--- Slide: MinMax with brace annotations ---
  \begin{frame}{\textit{MinMax Loss} -- Papéis de gerador e discriminador}
    \vspace{-0.3em}
    \[
      \hat{\phi},\hat{\thetav} =
      \argmin_{\phi}\, \argmax_{\thetav}
      \Bigg[\underbrace{ \sum_{j}
        \left(-\ln\!\left[1-\sigma\!\left(f\!\left[\textcolor{GenRed}{g[z^{(j)},\thetav]},\phi\right]\right)\right]\right)}_{\textcolor{GenRed}{\text{Gerador + Discriminador}}}
        \underbrace{- \sum_{i}
        \left(\ln\!\left[\sigma\!\left(f\!\left[x^{(i)},\phi\right]\right)\right]\right)}_{\textcolor{DiscBlue}{\text{Discriminador}}} \Bigg]
    \]
    \vspace{0.5em}
    \begin{itemize}
      \item O \textcolor{DiscBlue}{discriminador} recebe gradientes
      dos dois termos: quer classificar bem reais \emph{e} geradas.
      \item O \textcolor{GenRed}{gerador} só interage com o primeiro
      termo: só se importa em enganar o discriminador sobre as
      amostras sintéticas.
    \end{itemize}
  \end{frame}

  %--- Slide: Loss por rede ---
  \begin{frame}{\textit{MinMax Loss} -- Loss por rede}
    Na prática, treinamos as redes alternadamente, cada uma otimizando
    sua própria loss:
    \vspace{0.8em}
    \begin{block}{Loss do discriminador (minimizar em \textcolor{DiscBlue}{$\phi$})}
      \[
        L[\textcolor{DiscBlue}{\phi}] = \sum_{j}
          \!\left(-\ln\!\left[1-\sigma\!\left(f\!\left[g[z^{(j)},\thetav],\textcolor{DiscBlue}{\phi}\right]\right)\right]\right)
          - \sum_{i}
          \!\left(\ln\!\left[\sigma\!\left(f\!\left[x^{(i)},\textcolor{DiscBlue}{\phi}\right]\right)\right]\right)
      \]
    \end{block}
    \begin{block}{Loss do gerador (minimizar em \textcolor{GenRed}{$\thetav$})}
      \[
        L[\textcolor{GenRed}{\thetav}] = \sum_{j}
          \!\left(\ln\!\left[1-\sigma\!\left(f\!\left[g[z^{(j)},\textcolor{GenRed}{\thetav}],\phi\right]\right)\right]\right)
      \]
    \end{block}
  \end{frame}

  %--- Slide: Non-saturating loss ---
  \begin{frame}{\textit{Non-saturating Loss} -- Truque prático do gerador}
    \begin{columns}[T]
      \begin{column}{0.50\textwidth}
        No início do treino o discriminador vence facilmente:
        $\sigma(f[g(z),\phi]) \approx 0$. A loss $\ln[1-\sigma(\cdot)]$
        fica quase plana em $0$, e o gradiente que chega ao gerador é
        \emph{minúsculo}.
        \vspace{0.3em}
        \begin{block}{Substituição usada na prática}
          \[
            L_{\text{NS}}[\thetav] = \sum_{j}
              \!\left(-\ln\!\left[\sigma\!\left(f\!\left[g[z^{(j)},\thetav],\phi\right]\right)\right]\right)
          \]
        \end{block}
        \footnotesize
        \begin{itemize}
          \item Em vez de \emph{minimizar a probabilidade de ser pega}, o
          gerador \emph{maximiza a probabilidade de parecer real}.
          \item Mesmo ponto fixo (equilíbrio em $\sigma = \tfrac{1}{2}$),
          mas gradiente saudável quando o gerador é fraco.
          \item Já estava na proposta original de
          \citeauthoryear{goodfellow2014generative} e é o default em
          praticamente todas as implementações.
        \end{itemize}
      \end{column}
      \begin{column}{0.48\textwidth}
        \centering
        \begin{tikzpicture}[font=\small]
          \begin{axis}[
              width=6.2cm, height=5.2cm,
              axis lines=left,
              xmin=-6, xmax=6, ymin=0, ymax=1.05,
              xlabel={$d = f[g(z),\phi]$},
              ylabel={$|\partial L/\partial d|$},
              ytick={0,0.5,1},
              xtick={-6,-3,0,3,6},
              legend style={
                at={(0.5,-0.30)}, anchor=north,
                legend columns=1, font=\tiny,
                draw=none, fill=white,
              },
              clip=false,
            ]
            \addplot[GenRed, thick, dashed, samples=200, domain=-6:6]
              {1/(1+exp(-x))};
            \addlegendentry{original: $|\partial L/\partial d| = \sigma(d)$}
            \addplot[DiscBlue, thick, samples=200, domain=-6:6]
              {1 - 1/(1+exp(-x))};
            \addlegendentry{NS: $|\partial L/\partial d| = 1-\sigma(d)$}
            \draw[<-, gray] (axis cs:-5.0,0.03) -- (axis cs:-3.5,0.20)
              node[right, font=\scriptsize, text=gray] {grad $\approx 0$};
            \draw[<-, gray] (axis cs:-5.0,0.97) -- (axis cs:-3.5,0.80)
              node[right, font=\scriptsize, text=gray] {grad forte};
          \end{axis}
        \end{tikzpicture}\\
        \vspace{0.4em}
        \tiny Quando o discriminador vence ($\sigma \approx 0$,
        $d \to -\infty$) o gradiente da loss original some, o da NS
        permanece próximo de 1. O gráfico mostra $|\partial L/\partial d|$,
        os sinais reais são negativos ($-\sigma(d)$ e $\sigma(d)-1$).
      \end{column}
    \end{columns}
  \end{frame}

  %--- Slide: Algoritmo de treino ---
  \begin{frame}{\textit{MinMax Loss} -- Algoritmo de treino}
    Treinamos as redes \textbf{alternadamente}, em passos de gradiente:
    \vspace{0.3em}
    \begin{algorithmic}[1]
      \Require dataset real, número de passos do discriminador $k$,
               taxas de aprendizado $\eta_\phi$, $\eta_\theta$,
               tamanho de batch $m$
      \While{$\thetav$ não convergiu}
        \For{$t = 1, \ldots, k$}
          \State Amostre minibatch de reais $\{x^{(i)}\}_{i=1}^{m}$
          \State Amostre ruído $\{z^{(j)}\}_{j=1}^{m}$
          \State $\phi \leftarrow \phi - \eta_\phi \nabla_\phi L[\phi]$
        \EndFor
        \State Amostre novo ruído $\{z^{(j)}\}_{j=1}^{m}$
        \State $\thetav \leftarrow \thetav - \eta_\theta \nabla_\theta L[\thetav]$
      \EndWhile
    \end{algorithmic}
    \vspace{0.3em}
    \begin{itemize}
      \item No paper original $k=1$. Em \textit{WGAN} veremos $k=5$
      (queremos crítico próximo do ótimo).
      \item Otimizador típico: \textsc{Adam} com $\beta_1 = 0.5$
      ($\beta_1 = 0.9$ desestabiliza GANs).
    \end{itemize}
  \end{frame}

  %======================================================================
  % SECTION 3: DCGAN
  %======================================================================
  \section{DCGAN}

  %--- Slide: DCGAN architecture ---
  \begin{frame}{\textit{DCGAN}\footfullcite{radford2015dcgan}}
    A primeira arquitetura a popularizar GANs para imagens naturais:
    \textbf{Deep Convolutional GAN}.
    \vspace{0.3em}
    \centering
    \includegraphics[width=0.95\linewidth]{slides/img/05-geradores/dcgan-architecture.png}\\
    \vspace{0.4em}
    \begin{itemize}
      \item Gerador usa \textbf{convoluções fracionárias} (upsampling)
      partindo de um vetor latente $100 \times 1$.
      \item Discriminador usa \textbf{convoluções com stride} para reduzir
      a resolução.
      \item \emph{BatchNorm}, ReLU no gerador, LeakyReLU no discriminador.
    \end{itemize}
  \end{frame}

  %--- Slide: DCGAN guidelines ---
  \begin{frame}{\textit{DCGAN} -- Receita arquitetural}
    O paper destila um conjunto de \emph{guidelines} que dominaram os
    anos seguintes:
    \begin{itemize}
      \item \textbf{Sem \textit{pooling}}: troque por convoluções com
      \textit{stride} (D) e convoluções fracionárias (G).
      \item \textbf{Sem camadas \textit{fully connected}} ocultas em
      ambas as redes.
      \item \textbf{BatchNorm} em D e G, exceto na saída de G e na
      entrada de D (estabiliza o treino sem distorcer a saída).
      \item \textbf{Ativações}: ReLU em todas as camadas de G (Tanh na
      saída); LeakyReLU em todas as camadas de D.
      \item Vetor latente $z \in \mathbb{R}^{100}$, normal padrão.
    \end{itemize}
    \begin{block}{Por que importa}
      Antes da DCGAN, treinar GAN convolucional era considerado quase
      impossível. Esses cinco pontos viraram \emph{baseline} para
      praticamente todas as arquiteturas posteriores.
    \end{block}
  \end{frame}

  %--- Slide: DCGAN samples ---
  \begin{frame}{\textit{DCGAN} -- Amostras geradas}
    \centering
    \includegraphics[width=0.75\linewidth]{slides/img/05-geradores/dcgan-samples.png}\\
    \vspace{0.3em}
    \tiny Amostras sintéticas geradas por DCGANs treinadas em três
    \textit{datasets}: (a) faces, (b) ImageNet, (c) LSUN (cenas).
    Adaptado de \citeauthoryear{radford2015dcgan}.
    \note{
      As amostras já eram impressionantes para 2015. Hoje são facilmente
      superadas por \emph{StyleGAN} e por modelos difusores, mas a DCGAN
      estabeleceu o padrão arquitetural que dominou os anos seguintes.
    }
  \end{frame}

  %======================================================================
  % SECTION 4: COLAPSO PARA A MODA
  %======================================================================
  \section{Colapso para a Moda}

  %--- Slide: Mode collapse - descrição ---
  \begin{frame}{\textit{Mode Collapse}}
    É comum que o gerador \textbf{``colapse para a moda''} da
    distribuição, gerando instâncias sem diversidade.
    \vspace{0.5em}
    \begin{columns}
      \begin{column}{0.48\textwidth}
        \centering
        \includegraphics[width=\linewidth]{slides/img/05-geradores/mode-collapse-lsun.png}\\
        \tiny GAN em LSUN: cenas muito similares.
      \end{column}
      \begin{column}{0.48\textwidth}
        \centering
        \includegraphics[width=\linewidth]{slides/img/05-geradores/mode-collapse-mnist.png}\\
        \tiny MNIST ao longo das iterações: diversidade some.
      \end{column}
    \end{columns}
    \vspace{0.5em}
    O motivo disso fica aparente ao analisarmos a função de custo.
    \note{
      \emph{Mode collapse} é um dos fracassos mais famosos das GANs: o
      gerador encontra um conjunto pequeno de amostras que enganam bem
      o discriminador e fica preso ali.
    }
  \end{frame}

  %======================================================================
  % SECTION 5: ANÁLISE DA LOSS
  %======================================================================
  \section{Análise da Loss}

  %--- Slide: Divide by number ---
  \begin{frame}{Análise da Loss -- Normalização}
    Partimos da loss do discriminador, escrita aqui sem o sinal de menos
    (a forma que o discriminador \textbf{maximiza}), e \textbf{dividimos
    cada somatório} pelo número de instâncias (geradas $J$ ou reais $I$):
    \vspace{0.4em}
    \[
      L[\phi] = \sum_{j}
        \!\!\left(\ln\!\left[1-\sigma\!\left(f[x^{(j)*},\phi]\right)\right]\right)
        + \sum_{i} \!\!\left(\ln\!\left[\sigma\!\left(f[x^{(i)},\phi]\right)\right]\right)
    \]
    \[
      L[\phi] \approx \tfrac{1}{J}\sum_{j}
        \!\!\left(\ln\!\left[1-\sigma\!\left(f[x^{(j)*},\phi]\right)\right]\right)
        + \tfrac{1}{I}\sum_{i} \!\!\left(\ln\!\left[\sigma\!\left(f[x^{(i)},\phi]\right)\right]\right)
    \]
    \vspace{0.3em}
    Isso não altera o \emph{mínimo} (multiplicar por constante não
    muda argmin), mas prepara o próximo passo.
  \end{frame}

  %--- Slide: Expectation ---
  \begin{frame}{Análise da Loss -- Esperança/Expectativa}
    Ao dividir pelo número de instâncias, cada termo vira uma
    \textbf{esperança} sobre a distribuição correspondente:
    \vspace{0.4em}
    \[
      L[\phi] = \mathbb{E}_{x^{*}}\!\left[\ln\!\left[1-\sigma(f[x^{*},\phi])\right]\right]
      + \mathbb{E}_{x}\!\left[\ln\!\left[\sigma(f[x,\phi])\right]\right]
    \]
    \begin{itemize}
      \item $\mathbb{E}_{x^*}$: esperança sobre a distribuição gerada
      $P(x^*)$.
      \item $\mathbb{E}_{x}$: esperança sobre a distribuição real $P(x)$.
      \item Essa reescrita muda o ponto de vista: deixamos de falar de
      amostras e passamos a falar de distribuições.
    \end{itemize}
  \end{frame}

  %--- Slide: Integral ---
  \begin{frame}{Análise da Loss -- Integrais}
    Aplicando a definição de esperança, ficamos com \textbf{duas
    integrais}:
    \vspace{0.4em}
    \[
      L[\phi] = \int P(x^*) \ln\!\left[1-\sigma(f[x^{*},\phi])\right] dx^{*}
      + \int P(x) \ln\!\left[\sigma(f[x,\phi])\right] dx
    \]
    \vspace{0.3em}
    \begin{itemize}
      \item Os integrandos dependem do discriminador $\sigma(f[\cdot,\phi])$.
      \item Vamos substituir o discriminador pela expressão \emph{ótima}
      (que minimiza essa loss).
    \end{itemize}
  \end{frame}

  %--- Slide: Optimal discriminator ---
  \begin{frame}{Análise da Loss -- Discriminador ótimo}
    A saída $\sigma\!\left[f[\tilde{x},\phi]\right]$ é a probabilidade
    que o discriminador atribui ao rótulo \emph{real} dada uma amostra
    $\tilde{x}$. Por \textbf{regra de Bayes}:
    \[
      \sigma\!\left[f[\tilde{x},\phi]\right]
      = P(\text{real}\mid \tilde{x})
      = \frac{P(\tilde{x}\mid \text{real})\,P(\text{real})}
             {P(\tilde{x}\mid \text{real})\,P(\text{real})
              + P(\tilde{x}\mid \text{gen})\,P(\text{gen})}
    \]
    Assumindo $I = J$, ou seja $P(\text{real}) = P(\text{gen}) = \tfrac{1}{2}$,
    os \emph{priors} se cancelam:
    \[
      \sigma\!\left[f[\tilde{x},\phi]\right]
      = \frac{P(\tilde{x}\mid \text{real})}
             {P(\tilde{x}\mid \text{real}) + P(\tilde{x}\mid \text{gen})}
      = \frac{P(x)}{P(x)+P(x^{*})}
    \]
    \footnotesize
    \begin{itemize}
      \item $P(x)$ é a densidade dos dados reais em $\tilde{x}$, $P(x^{*})$
      é a densidade gerada no mesmo ponto.
      \item Em palavras: a probabilidade ideal de ser real é a fração da
      densidade total naquele ponto que vem do conjunto real.
      \item É o classificador perfeito, depende apenas das densidades, não
      de uma rede específica.
    \end{itemize}
    \note{
      Para provar: tome o integrando, derive em relação a $\sigma$ e
      iguale a zero. A solução é essa razão.
    }
  \end{frame}

  %--- Slide: Substituting ---
  \begin{frame}{Análise da Loss -- Substituindo o discriminador ótimo}
    Substituindo $\sigma(f[x,\phi]) = \frac{P(x)}{P(x)+P(x^*)}$ na loss:
    \vspace{0.2em}
    \[
      L[\phi] = \int P(x^{*}) \ln\!\left[\frac{P(x^{*})}{P(x)+P(x^{*})}\right] dx^{*}
      + \int P(x)\ln\!\left[\frac{P(x)}{P(x)+P(x^{*})}\right] dx
    \]
    \vspace{0.3em}
    Os dois termos dentro dos logaritmos viraram \textbf{razões de
    densidades}, que é exatamente a forma de uma divergência KL.
    \note{
      Note que o primeiro termo começou como $\ln[1-\sigma]$. Como
      $1 - \tfrac{P(x)}{P(x)+P(x^*)} = \tfrac{P(x^*)}{P(x)+P(x^*)}$,
      chegamos à forma acima.
    }
  \end{frame}

  %--- Slide: Multiply and divide by 2 ---
  \begin{frame}{Análise da Loss -- Truque do fator 2}
    Por conveniência matemática, multiplicamos e dividimos por 2 dentro
    do logaritmo (não altera o $\argmin$):
    \vspace{0.2em}
    \[
      L[\phi] = \int P(x^{*}) \ln\!\left[\frac{2P(x^{*})}{P(x)+P(x^{*})}\right] dx^{*}
      + \int P(x)\ln\!\left[\frac{2P(x)}{P(x)+P(x^{*})}\right] dx - 2\ln 2
    \]
    \vspace{0.3em}
    Os dois fatores $\ln 2$ saem das integrais (pois $\int P = \int P^{*} = 1$).
    Agora cada termo é exatamente uma divergência KL entre uma
    distribuição e a \emph{mistura} $\frac{P(x^*)+P(x)}{2}$.
  \end{frame}

  %--- Slide: Recognizing KL ---
  \begin{frame}{Análise da Loss -- Reconhecendo KL}
    Cada parcela tem a forma da \textbf{KL Divergence}:
    \vspace{0.3em}
    \[
      D_{KL}\!\left[p(x) \,\|\, q(x)\right] = \int p(x)\ln\!\left[\frac{p(x)}{q(x)}\right] dx
    \]
    \vspace{0.3em}
    Portanto:
    \[
      L[\phi] = D_{KL}\!\left[P(x^{*}) \,\Big\|\, \tfrac{P(x^{*})+P(x)}{2}\right]
             + D_{KL}\!\left[P(x) \,\Big\|\, \tfrac{P(x^{*})+P(x)}{2}\right] - 2\ln 2
    \]
  \end{frame}

  %--- Slide: Jensen-Shannon ---
  \begin{frame}{Análise da Loss -- \textit{Jensen--Shannon Divergence}}
    Reconhecemos a definição da \textbf{Jensen--Shannon Divergence}:
    \vspace{0.3em}
    \[
      D_{JS}\!\left[P(x^{*}) \,\|\, P(x)\right]
      = \tfrac{1}{2} D_{KL}\!\left[P(x^{*}) \,\Big\|\, \tfrac{P(x^{*})+P(x)}{2}\right]
      + \tfrac{1}{2} D_{KL}\!\left[P(x) \,\Big\|\, \tfrac{P(x^{*})+P(x)}{2}\right]
    \]
    \vspace{0.2em}
    Logo $L[\phi] = 2\,D_{JS}\!\left[P(x^{*}) \,\|\, P(x)\right] - 2\ln 2$.
    \begin{alertblock}{Interpretação}
      Com discriminador ótimo, a loss de uma GAN é, a menos de uma
      constante e de um fator, a \textbf{divergência de Jensen--Shannon}
      entre a distribuição real e a gerada.
    \end{alertblock}
  \end{frame}

  %--- Slide: qualidade x cobertura ---
  \begin{frame}{Análise da Loss -- Qualidade vs.\ Cobertura}
    \begin{columns}[T]
      \begin{column}{0.55\textwidth}
        \vspace{-0.4em}
        \[
          \resizebox{\linewidth}{!}{$
            D_{JS}[P(x^*)\|P(x)] =
            \underbrace{\tfrac{1}{2}D_{KL}\!\!\left[P(x^*)\Big\|\tfrac{P(x^*)+P(x)}{2}\right]}_{\textcolor{GenRed}{\text{Qualidade}}}
            + \underbrace{\tfrac{1}{2}D_{KL}\!\!\left[P(x)\Big\|\tfrac{P(x^*)+P(x)}{2}\right]}_{\textcolor{DiscBlue}{\text{Cobertura}}}
          $}
        \]
        \footnotesize
        \begin{itemize}
          \item \textcolor{GenRed}{\textbf{Primeiro termo (qualidade)}},
          penaliza regiões em que $P(x^*)$ tem probabilidade alta mas
          a mistura não tem, isto é, amostras geradas implausíveis.
          \item \textcolor{DiscBlue}{\textbf{Segundo termo (cobertura)}},
          penaliza regiões em que $P(x)$ tem probabilidade alta mas a
          mistura não tem, isto é, modas ignoradas pelo gerador.
        \end{itemize}
        \begin{block}{Pergunta}
          Qual dos dois termos o gerador consegue \emph{de fato}
          controlar?
        \end{block}
      \end{column}
      \begin{column}{0.43\textwidth}
        \centering
        % P(x) = 0.5*N(-1.5, 0.2) + 0.5*N(4, 0.2),  P*(x) = N(-1, 0.25).
        % Coeficientes: 0.5/sqrt(0.4*pi) = 0.4460; 1/sqrt(0.5*pi) = 0.7979.
        \begin{tikzpicture}[font=\small]
          \begin{axis}[
              width=6.2cm, height=3.2cm,
              axis lines=left,
              xmin=-4, xmax=7, ymin=0, ymax=0.92,
              xtick=\empty, ytick=\empty,
              legend style={
                at={(0.98,0.98)}, anchor=north east,
                legend columns=1, row sep=-0.25em,
                font=\tiny, draw=none,
                fill=white, fill opacity=0.75, text opacity=1,
              },
              clip=false,
            ]
            \addplot[DiscBlue, thick, samples=240, domain=-4:7]
              {0.4460*exp(-(x+1.5)^2/0.4) + 0.4460*exp(-(x-4)^2/0.4)};
            \addlegendentry{$P(x)$ real}
            \addplot[GenRed, thick, samples=240, domain=-4:7]
              {0.7979*exp(-(x+1)^2/0.5)};
            \addlegendentry{$P^{*}(x)$ gerada}
            \addplot[gray, thick, dashed, samples=240, domain=-4:7]
              {0.5*(0.4460*exp(-(x+1.5)^2/0.4) + 0.4460*exp(-(x-4)^2/0.4) + 0.7979*exp(-(x+1)^2/0.5))};
            \addlegendentry{$M = \tfrac{1}{2}(P{+}P^{*})$}
          \end{axis}
        \end{tikzpicture}\\[-0.1em]
        \begin{tikzpicture}[font=\small]
          \begin{axis}[
              width=6.2cm, height=3.4cm,
              axis x line=middle, axis y line=none,
              xmin=-4, xmax=7, ymin=-0.12, ymax=0.22,
              xtick=\empty,
              clip=false,
            ]
            % qualidade: (1/2) p* ln(p*/M), com sinal preservado.
            \addplot[GenRed, thick, samples=300, domain=-4:7,
                     fill=GenRed!22, draw=GenRed]
              {0.5 * (0.7979*exp(-(x+1)^2/0.5))
                * ln( (0.7979*exp(-(x+1)^2/0.5))
                     / (0.5*(0.4460*exp(-(x+1.5)^2/0.4)
                            + 0.4460*exp(-(x-4)^2/0.4)
                            + 0.7979*exp(-(x+1)^2/0.5))) )} \closedcycle;
            % cobertura: (1/2) p ln(p/M), com sinal preservado.
            \addplot[DiscBlue, thick, samples=300, domain=-4:7,
                     fill=DiscBlue!22, draw=DiscBlue]
              {0.5 * (0.4460*exp(-(x+1.5)^2/0.4) + 0.4460*exp(-(x-4)^2/0.4))
                * ln( (0.4460*exp(-(x+1.5)^2/0.4) + 0.4460*exp(-(x-4)^2/0.4))
                     / (0.5*(0.4460*exp(-(x+1.5)^2/0.4)
                            + 0.4460*exp(-(x-4)^2/0.4)
                            + 0.7979*exp(-(x+1)^2/0.5))) )} \closedcycle;
            \node[GenRed, font=\scriptsize, anchor=south]
              at (axis cs:-0.7,0.17) {qualidade};
            \node[DiscBlue, font=\scriptsize, anchor=south]
              at (axis cs:4,0.16) {cobertura};
          \end{axis}
        \end{tikzpicture}\\[0.1em]
        \tiny Acima: densidades normalizadas. Abaixo: integrando de
        $\tfrac{1}{2}D_{KL}[P^{*}\|M]$ (vermelho, qualidade) e
        $\tfrac{1}{2}D_{KL}[P\|M]$ (azul, cobertura). Áreas acima do
        eixo somam, abaixo subtraem; a integral líquida é o termo
        correspondente da $D_{JS}$. A moda de $P$ em $x \approx 4$,
        ignorada por $P^{*}$, domina a cobertura.
      \end{column}
    \end{columns}
  \end{frame}

  %--- Slide: segundo termo não depende de gerador ---
  \begin{frame}{Análise da Loss -- Por que há colapso de moda}
    \[
      D_{JS}[P(x^*)\|P(x)] =
      \tfrac{1}{2}D_{KL}\!\!\left[P(x^*)\Big\|\tfrac{P(x^*)+P(x)}{2}\right]
      + \tfrac{1}{2}\underbrace{D_{KL}\!\!\left[P(x)\Big\|\tfrac{P(x^*)+P(x)}{2}\right]}_{\text{depende fracamente de } P(x^*)}
    \]
    \vspace{0.4em}
    \begin{alertblock}{Consequência}
      Quando $P(x^*) \ll P(x)$ em uma região,
      $(P(x^*){+}P(x))/2 \approx P(x)/2$, então o argumento do log no
      \textcolor{DiscBlue}{segundo termo} satura em
      $\tfrac{P(x)}{P(x)/2} = 2$, e o integrando reduz-se a
      $\tfrac{1}{2}P(x)\log 2$, \emph{independente} de $P(x^*)$.
      \vspace{0.3em}

      O gerador é \textbf{incentivado muito mais fortemente a
      qualidade} que a cobertura $\to$ \textbf{colapso para a moda}.
    \end{alertblock}
  \end{frame}

  %--- Slide: Recap analítico ---
  \begin{frame}{Análise da Loss -- Recapitulando}
    A cadeia de manipulações que aplicamos:
    \begin{enumerate}
      \item Loss original = entropia cruzada binária do discriminador.
      \item Normalizamos cada soma $\to$ esperanças sobre $P(x)$ e $P(x^{*})$.
      \item Esperanças $\to$ integrais.
      \item Substituímos $\sigma(f)$ pelo discriminador
      \emph{ótimo} $\frac{P(x)}{P(x)+P(x^{*})}$.
      \item Multiplicamos e dividimos por $2$ dentro do logaritmo.
      \item Cada termo virou uma KL contra a mistura
      $\frac{P(x)+P(x^{*})}{2}$.
      \item Somando, obtemos $L[\phi] = 2\,D_{JS}\!\left[P(x^{*})\,\|\,P(x)\right] - 2\ln 2$.
    \end{enumerate}
    \begin{alertblock}{Moral}
      Treinar uma GAN com discriminador ótimo equivale a minimizar
      Jensen--Shannon. As fragilidades dessa divergência (assimetria
      qualidade vs cobertura, comportamento em suportes disjuntos)
      viram fragilidades do treinamento.
    \end{alertblock}
  \end{frame}

  %--- Slide: Exemplo numérico - o que cada valor representa ---
  \begin{frame}{Exemplo Numérico}{O Que Cada Valor Representa}
    Antes de calcular o $D_{JS}$, vamos ancorar cada símbolo em algo
    \textbf{concreto}. Imagine que estamos treinando uma GAN para gerar
    imagens em três categorias: \emph{gato}, \emph{cão} e
    \emph{pássaro}. Identifique cada uma com $a$, $b$, $c$.
    \vspace{0.3em}
    \begin{itemize}
      \item $P(x)$: fração das imagens \textbf{reais} (no dataset) em cada
            categoria. É fixo, a verdade.
      \item $P^{*}(x)$: fração das imagens \textbf{geradas} pelo modelo
            atual em cada categoria. Muda conforme o gerador treina.
      \item $D^{*}(x)$: probabilidade que o \textbf{discriminador ótimo}
            atribui à categoria $x$ de ser real. Onde $P$ domina, $D^{*}$
            tende a $1$; onde $P^{*}$ domina, $D^{*}$ tende a $0$.
      \item $M(x) = \tfrac{1}{2}(P + P^{*})$: distribuição de
            \textbf{mistura}, o referencial neutro contra o qual ambas as
            divergências KL são medidas.
      \item $D_{JS}(P^{*}\|P)$: o \textbf{único número} que resume a
            distância entre o que o gerador faz hoje e o que deveria
            fazer. Zero $\Rightarrow$ gerador perfeito;
            $\ln 2 \Rightarrow$ suportes disjuntos.
    \end{itemize}
    \note{
      Vale enfatizar: $P$ vem dos dados, $P^{*}$ vem do modelo. O
      discriminador é uma \emph{ferramenta} que extrai o sinal dessa
      diferença e empurra o gerador via gradiente.
    }
  \end{frame}

  %--- Slide: Exemplo numérico JS divergence ---
  \begin{frame}{Exemplo Numérico: $D^{*}$ e $D_{JS}$ em 1D}
    Suporte $\{a, b, c\}$.
    Real $P = (0.5,\, 0.3,\, 0.2)$;\quad
    gerada $P^{*} = (0.2,\, 0.5,\, 0.3)$.
    \begin{columns}[T]
      \begin{column}{0.46\textwidth}
        \textbf{Discriminador ótimo} $D^{*}(x) = \frac{P(x)}{P(x) + P^{*}(x)}$:
        \[
          D^{*} \approx (\, 0.71,\; 0.38,\; 0.40 \,)
        \]
        \textbf{Mistura} $M = \tfrac{1}{2}(P + P^{*})$:
        \[
          M = (\, 0.35,\; 0.40,\; 0.25 \,)
        \]
      \end{column}
      \begin{column}{0.52\textwidth}
        \centering
        \resizebox{\linewidth}{!}{%
        \begin{tabular}{@{}c|c|c@{}}
          \toprule
          $x$      & $P\,\ln(P/M)$  & $P^{*}\,\ln(P^{*}/M)$ \\
          \midrule
          $a$      & $+0.178$     & $-0.112$            \\
          $b$      & $-0.086$     & $+0.112$            \\
          $c$      & $-0.045$     & $+0.055$            \\
          \midrule
          $\Sigma$ & $D_{KL}(P\|M) = 0.047$ & $D_{KL}(P^{*}\|M) = 0.054$ \\
          \bottomrule
        \end{tabular}}
      \end{column}
    \end{columns}
    \[
      D_{JS}(P^{*} \| P) = \tfrac{1}{2}(0.047 + 0.054) \approx \mathbf{0.05}\ \text{nats}
    \]
    \note{
      Em bits (log base 2), o máximo de $D_{JS}$ é $1$; em nats,
      $\ln 2 \approx 0.693$. Diferente do $D_{KL}$, $D_{JS}$ é sempre
      limitado, mas ainda satura quando os suportes não se encontram,
      o que produz gradientes fracos para o gerador.
    }
  \end{frame}

  %--- Slide: Exemplo numérico - limite de suporte disjunto ---
  \begin{frame}{Exemplo Numérico: limite de suporte disjunto}
    \footnotesize
    Mesmo suporte $\{a, b, c\}$, mas agora a categoria $c$ \textbf{não
    aparece} nos dados reais: $P = (0.6,\, 0.4,\, 0.0)$, gerada
    $P^{*} = (0.2,\, 0.5,\, 0.3)$.
    Mistura $M = (\, 0.40,\; 0.45,\; 0.15 \,)$.
    \begin{columns}[T]
      \begin{column}{0.46\textwidth}
        Em $x = c$ temos $P(c) = 0$: o termo $P\ln(P/M)$ se anula
        (convenção $0\ln 0 = 0$). Já $P^{*}(c) = 0.3$ e
        $M(c) = P^{*}(c)/2$, então
        \[
          P^{*}(c)\,\ln\!\frac{P^{*}(c)}{M(c)} = 0.3 \cdot \ln 2 \approx 0.21.
        \]
      \end{column}
      \begin{column}{0.52\textwidth}
        \centering
        \resizebox{\linewidth}{!}{%
        \begin{tabular}{@{}c|c|c@{}}
          \toprule
          $x$      & $P\,\ln(P/M)$ & $P^{*}\,\ln(P^{*}/M)$ \\
          \midrule
          $a$      & $+0.243$    & $-0.139$            \\
          $b$      & $-0.047$    & $+0.053$            \\
          $c$      & $\;\;\;0$     & $+0.208$            \\
          \midrule
          $\Sigma$ & $D_{KL}(P\|M)\!=\!0.196$ & $D_{KL}(P^{*}\|M)\!=\!0.122$ \\
          \bottomrule
        \end{tabular}}
        \[
          D_{JS}(P^{*} \| P) \approx \mathbf{0.16}\ \text{nats}
        \]
      \end{column}
    \end{columns}
    \begin{alertblock}{Limite de suporte totalmente disjunto}
      Quando $P$ e $P^{*}$ não compartilham nenhum ponto de suporte,
      $D_{JS}(P^{*} \| P) \to \ln 2 \approx 0.69$ nats (máximo da
      Jensen--Shannon). $D_{JS}$ trava nesse valor e o gradiente do
      gerador desaparece, motivando a distância de Wasserstein.
    \end{alertblock}
  \end{frame}

  %======================================================================
  % SECTION 6: VANISHING GRADIENT
  %======================================================================
  \section{Vanishing Gradient}

  %--- Slide: Manifold hypothesis ---
  \begin{frame}{Por que distribuições disjuntas são a regra}
    \begin{itemize}
      \item Dados naturais ocupam apenas uma pequena fatia do espaço de
      pixels: rostos, por exemplo, formam uma superfície curva de
      dimensão muito menor que a dimensão da imagem. Chamamos essa
      superfície de \emph{manifold}.
      \item O gerador, por construção, também produz amostras sobre um
      manifold: a imagem da função $g[\cdot,\thetav]$ aplicada ao
      latente $z \in \mathbb{R}^{d}$, com $d \ll D$.
      \item Dois manifolds de dimensão baixa em $\mathbb{R}^{D}$,
      escolhidos independentemente, quase nunca se cruzam: a interseção
      tem \textbf{medida zero}.
      \item Consequência: no início do treino, $P(x)$ e $P(x^{*})$ vivem
      em regiões praticamente disjuntas, e o $D_{JS}$ fica preso em
      $\ln 2$.
    \end{itemize}
    \begin{alertblock}{Conclusão}
      Distribuições disjuntas \textbf{não} são caso patológico raro:
      são a situação típica para dados de alta dimensão. Por isso o
      problema a seguir aparece em quase todo treino de GAN.
    \end{alertblock}
  \end{frame}

  %--- Slide: Vanishing gradient - teoria ---
  \begin{frame}{\textit{Vanishing Gradient} -- Motivação}
    \begin{itemize}
      \item Vimos que, no caso de discriminador ótimo, a loss minimiza
      uma \textbf{distância} ($D_{JS}$) entre instâncias reais e
      geradas.
      \item Contudo, essa abordagem traz problemas:
      \begin{itemize}
        \item Se as distribuições são \textbf{disjuntas}, a $D_{JS}$
        assume valor constante $\ln 2$ e a $D_{KL}$ diverge.
        \item Nesse caso, \emph{nenhuma ``pequena'' mudança} no gerador
        melhora a loss, os gradientes desaparecem.
      \end{itemize}
    \end{itemize}
    \begin{block}{Relembrando KL}
      \[
        D_{KL}[p(x)\|q(x)] = \int p(x)\ln\!\left[\frac{p(x)}{q(x)}\right] dx
      \]
      Diverge sempre que existe $x$ com $p(x)>0$ e $q(x)=0$.
    \end{block}
  \end{frame}

  %--- Slide: Vanishing gradient - empírico ---
  \begin{frame}{\textit{Vanishing Gradient} -- Evidência empírica\footfullcite{arjovsky2017towards}}
    \begin{columns}[T]
      \begin{column}{0.48\textwidth}
        \begin{itemize}
          \item O discriminador é congelado após $1$, $10$ e $25$ épocas
          e o gerador continua treinando.
          \item Conforme o discriminador fica ``bom demais'', o
          gradiente que chega ao gerador despenca (escala log).
          \item Sintoma típico: treinamento \emph{instável}, no qual o
          discriminador precisa ser \emph{deliberadamente enfraquecido}
          para que o gerador aprenda.
        \end{itemize}
      \end{column}
      \begin{column}{0.50\textwidth}
        \centering
        \includegraphics[width=\linewidth]{slides/img/05-geradores/vanishing-gradient.png}\\
        \tiny Adaptado de \citeauthoryear{arjovsky2017towards}.
      \end{column}
    \end{columns}
  \end{frame}

  %--- Slide: Pedindo outra distância ---
  \begin{frame}{Motivação para mudar a distância}
    \begin{itemize}
      \item Vimos que a loss original de GANs pode ser interpretada
      como uma \textbf{distância entre distribuições de probabilidade}.
      \item Porém, é uma distância com propriedades \textbf{inadequadas
      para otimização}:
      \begin{itemize}
        \item Distribuições disjuntas têm $D_{KL}$ infinita e
        $D_{JS}$ constante.
        \item Gradientes somem rapidamente.
      \end{itemize}
      \item Vamos escolher uma distância com propriedades mais
      adequadas para otimização: a distância de \textbf{Wasserstein}.
    \end{itemize}
  \end{frame}

  %======================================================================
  % SECTION 7: WASSERSTEIN GAN
  %======================================================================
  \section{Wasserstein GAN}

  %--- Slide: Wasserstein Distance definição ---
  \begin{frame}{\textit{Wasserstein Distance}\footfullcite{arjovsky2017wasserstein}}
    \begin{itemize}
      \item Também conhecida como \textit{Earth Mover's Distance} (para
      distribuições discretas).
      \item Definida pela \textbf{quantidade de trabalho} necessária
      para transportar a massa de probabilidade de uma distribuição até
      transformá-la em outra.
      \begin{itemize}
        \item \emph{Trabalho} = massa $\times$ distância movida.
      \end{itemize}
      \item Oferece \textbf{gradientes informativos} mesmo quando as
      distribuições são disjuntas.
    \end{itemize}
  \end{frame}

  %--- Slide: Wasserstein discreta ---
  \begin{frame}{\textit{Wasserstein Distance} -- Caso discreto}
    \begin{block}{Definição}
      \[
        D_{w}\!\left[P(x)\,\|\,Q(x)\right] =
        \min_{T}\!\left[\sum_{i,j} T_{ij}\,|i-j|\right]
      \]
    \end{block}
    \begin{itemize}
      \item $T$ é o \textbf{``plano de transporte''}: $T_{ij}$ indica
      quanta massa vai de $i$ (origem) para $j$ (destino).
      \item Sujeito às restrições:
      \begin{itemize}
        \item $\sum_{j} T_{ij} = P(x=i)$
        \item $\sum_{i} T_{ij} = Q(x=j)$
        \item $T_{ij} \geq 0$
      \end{itemize}
    \end{itemize}
  \end{frame}

  %--- Slide: Transporte visual ---
  \begin{frame}{\textit{Wasserstein Distance} -- Visualização}
    \centering
    \begin{tikzpicture}
      \definecolor{otlow}{RGB}{42,13,10}
      \definecolor{othigh}{RGB}{255,236,190}
      \def\s{0.42}
      % ---------- Distribuição P ----------
      \begin{scope}[shift={(0,2.7)}]
        \foreach \i/\v in {0/0.10,1/0.15,2/0.05,3/0.05,4/0.20,5/0.25,6/0.15,7/0.05}
          \fill[gray!65] ({\i*\s+0.03},0) rectangle ({(\i+1)*\s-0.03},{\v*6.2});
        \draw[gray!60] (0,0) -- ({8*\s},0);
        \node[font=\scriptsize, anchor=south] at ({4*\s},1.6) {$P(x{=}i)$};
        \node[font=\tiny, anchor=north, gray] at ({4*\s},-0.04) {índice $i$};
      \end{scope}
      % ---------- Distribuição Q ----------
      \begin{scope}[shift={(0,-0.4)}]
        \foreach \j/\v in {0/0.05,1/0.10,2/0.20,3/0.25,4/0.20,5/0.10,6/0.05,7/0.05}
          \fill[teal!55] ({\j*\s+0.03},0) rectangle ({(\j+1)*\s-0.03},{\v*6.2});
        \draw[gray!60] (0,0) -- ({8*\s},0);
        \node[font=\scriptsize, anchor=south] at ({4*\s},1.6) {$Q(x{=}j)$};
        \node[font=\tiny, anchor=north, gray] at ({4*\s},-0.04) {índice $j$};
      \end{scope}
      % ---------- Plano de transporte T (acoplamento monótono ótimo) ----------
      \begin{scope}[shift={(5.2,0.3)}]
        \fill[otlow] (0,0) rectangle ({8*\s},{8*\s});
        \foreach \i/\j/\v in {0/0/0.05,0/1/0.05,1/1/0.05,1/2/0.10,2/2/0.05,
                              3/2/0.05,4/3/0.20,5/3/0.05,5/4/0.20,6/5/0.10,
                              6/6/0.05,7/7/0.05}{
          \pgfmathsetmacro\f{\v/0.20*100}
          \fill[othigh!\f!otlow] ({\j*\s},{(7-\i)*\s}) rectangle ++(\s,\s);
        }
        \draw[gray!60, thin] (0,0) rectangle ({8*\s},{8*\s});
        \draw[cyan, very thick] ({4*\s},{2*\s}) rectangle ++(\s,\s);
        % marginais: P à esquerda, Q abaixo
        \foreach \i/\v in {0/0.10,1/0.15,2/0.05,3/0.05,4/0.20,5/0.25,6/0.15,7/0.05}
          \fill[gray!65] (-0.08,{(7-\i)*\s+0.03}) rectangle ({-0.08-\v*2.4},{(8-\i)*\s-0.03});
        \foreach \j/\v in {0/0.05,1/0.10,2/0.20,3/0.25,4/0.20,5/0.10,6/0.05,7/0.05}
          \fill[teal!55] ({\j*\s+0.03},-0.08) rectangle ({(\j+1)*\s-0.03},{-0.08-\v*2.4});
        \node[font=\tiny, anchor=east, gray] at ({-0.08-0.25*2.4},{2.5*\s}) {$i{=}5$};
        \node[font=\tiny, anchor=north, gray] at ({4.5*\s},{-0.08-0.20*2.4}) {$j{=}4$};
        \node[font=\scriptsize, anchor=south] at ({4*\s},{8*\s+0.08}) {Plano de transporte $T$};
      \end{scope}
      % ---------- Matriz de custo |i-j| ----------
      \begin{scope}[shift={(10.0,0.3)}]
        \foreach \i in {0,...,7}
          \foreach \j in {0,...,7}{
            \pgfmathsetmacro\f{abs(\i-\j)/7*100}
            \fill[othigh!\f!otlow] ({\j*\s},{(7-\i)*\s}) rectangle ++(\s,\s);
          }
        \draw[gray!60, thin] (0,0) rectangle ({8*\s},{8*\s});
        \node[font=\scriptsize, anchor=south] at ({4*\s},{8*\s+0.08}) {Custo $|i-j|$};
        \node[font=\scriptsize, anchor=north] at ({4*\s},-0.35) {$D_{w} = \sum_{i,j} T_{ij}\,|i-j|$};
      \end{scope}
    \end{tikzpicture}\\
    \vspace{0.3em}
    \tiny Exemplo discreto com suporte $\{0,\dots,7\}$: as marginais do
    plano $T$ (barras à esquerda e abaixo da matriz) reproduzem $P$ e $Q$,
    e a célula destacada é $T_{54} = 0.2$. Neste exemplo, $D_{w} = 0.8$.
    \note{
      A matriz de transporte só concentra massa na diagonal quando as
      distribuições já são iguais; quanto mais precisa ``empurrar''
      massa, mais a matriz se afasta da diagonal.
    }
  \end{frame}

  %--- Slide: Primal vs Dual ---
  \begin{frame}{\textit{Wasserstein Distance} -- Formulação dual}
    \begin{itemize}
      \item Para obter o valor de $D_{w}$, precisamos resolver um
      problema de \textbf{programação linear} em $T$.
      \item Isso não é prático durante o treinamento de GANs.
      \item Para esse problema específico, a formulação \textbf{dual} é
      mais útil.
    \end{itemize}
  \end{frame}

  %--- Slide: De onde vem o f (LP duality) ---
  \begin{frame}{\textit{Wasserstein Distance} -- De onde vem o $f$?}
    \begin{itemize}
      \item O primal do slide anterior é um \textbf{LP}: minimizamos
      $\sum_{i,j} T_{ij}\,|i-j|$ sobre o plano $T$, sujeito às marginais
      e a $T_{ij}\ge 0$.
      \item Todo LP tem um \textbf{dual}, com uma variável por restrição.
      As marginais geram as variáveis $f_i$ (nas origens) e $g_j$ (nos
      destinos), e o dual \emph{maximiza}:
      \[
        \max_{f,g}\;\sum_i P(x{=}i)\,f_i + \sum_j Q(x{=}j)\,g_j
        \quad\text{s.a.}\quad f_i + g_j \le |i-j|.
      \]
      \item \emph{Kantorovich--Rubinstein}: para o custo $|i-j|$, o ótimo
      ocorre com $g=-f$. A restrição vira $f_i - f_j \le |i-j|$, que é
      exatamente $f$ ser \textbf{1-Lipschitz}.
      \item Sobra uma única função $f$ a maximizar (próximo slide). É daí
      que vem o $f$: a variável dual associada às marginais.
    \end{itemize}
    \note{
      Não precisamos derivar o dual de cor; a mensagem é que $f$ não cai
      do céu, é a variável dual das restrições de marginal, e a condição
      $g=-f$ com $f_i-f_j\le|i-j|$ é o que chamamos de 1-Lipschitz.
    }
  \end{frame}

  %--- Slide: Dual form ---
  \begin{frame}{\textit{Wasserstein Distance} -- Forma dual}
    Reescrevendo em sua forma dual, temos uma nova variável de
    otimização $f$ que \textbf{maximizamos} sob a restrição de que
    valores adjacentes não se alterem em mais de uma unidade:
    \vspace{0.3em}
    \[
      D_{w}\!\left[P(x)\,\|\,Q(x)\right] =
      \max_{f}\!\left[\sum_{i} P(x=i)\,f_{i} - \sum_{j} Q(x=j)\,f_{j}\right]
    \]
    \vspace{0.3em}
    Sob a restrição: $|f_{i+1} - f_{i}| \le 1$
    (restrição de \textbf{1-Lipschitz} no caso contínuo).
    \note{
      O dual aparece por \emph{Kantorovich--Rubinstein duality}. A
      restrição de Lipschitz 1 torna o problema bem definido.
    }
  \end{frame}

  %--- Slide: Wasserstein loss no contínuo ---
  \begin{frame}{\textit{Wasserstein GAN} -- Loss}
    Considerando que as distribuições são aproximadas pelas instâncias
    reais e geradas, obtemos a loss do \textbf{WGAN}:
    \vspace{0.3em}
    \[
      L[\phi,\thetav] = \sum_{j} f\!\left[g[z^{(j)},\thetav],\phi\right]
      - \sum_{i} f\!\left[x^{(i)},\phi\right]
    \]
    \vspace{0.3em}
    Sujeito à restrição: $\left|\frac{\partial f(x,\phi)}{\partial x}\right| \le 1$
    \vspace{0.3em}
    \begin{itemize}
      \item Não há mais sigmóide: $f$ é um \textbf{crítico}, não um
      classificador.
      \item A restrição de Lipschitz é imposta via \emph{weight clipping}
      (a seguir) ou \emph{gradient penalty}\cite{gulrajani2017improved}.
    \end{itemize}
  \end{frame}

  %--- Slide: Weight clipping ---
  \begin{frame}{\textit{Wasserstein GAN} -- Weight clipping}
    \begin{itemize}
      \item Como garantir, na prática, que o crítico $f$ seja
      \textbf{1-Lipschitz} ao longo do treino?
      \item \emph{Weight clipping}: depois de cada atualização do
      crítico, fazemos o \emph{clip} de todos os pesos para o intervalo
      $[-c,\,c]$:
      \[
        w \leftarrow \operatorname{clip}(w, -c, c).
      \]
      \item Intuição: limitar a magnitude dos pesos limita o quão rápido
      $f$ pode variar, ou seja, limita a sua constante de Lipschitz. Não
      garante exatamente $1$, mas prende $f$ numa família de funções com
      inclinação controlada.
    \end{itemize}
    \begin{alertblock}{É um instrumento grosseiro}
      \begin{itemize}
        \item $c$ pequeno demais: os gradientes somem; $c$ grande demais:
        o treino fica instável.
        \item Os pesos tendem a saturar em $\pm c$, desperdiçando a
        capacidade do crítico (o que motiva o \emph{gradient penalty},
        adiante).
        \item Não confundir com \emph{gradient clipping}, que limita os
        gradientes, não os pesos.
      \end{itemize}
    \end{alertblock}
  \end{frame}

  %--- Slide: WGAN gradient comparison ---
  \begin{frame}{\textit{WGAN} -- Gradientes em distribuições disjuntas}
    \begin{columns}[T]
      \begin{column}{0.55\textwidth}
        \centering
        \includegraphics[width=\linewidth]{slides/img/05-geradores/wgan-gradients.png}\\
        \vspace{0.2em}
        \tiny Adaptado de Jonathan Hui (Medium, 2018).
      \end{column}
      \begin{column}{0.42\textwidth}
        \begin{itemize}
          \item Discriminador \textcolor{red}{saturado}
          $\Rightarrow$ gradientes nulos.
          \item Crítico WGAN produz gradientes \textbf{lineares} mesmo
          em regiões disjuntas $\Rightarrow$ treinamento muito mais
          estável.
          \item A restrição de Lipschitz é a chave: evita que o crítico
          ``exploda'' para separar as distribuições.
        \end{itemize}
      \end{column}
    \end{columns}
  \end{frame}

  %--- Slide: WGAN algorithm ---
  \begin{frame}{\textit{WGAN} -- Algoritmo}
    \begin{algorithmic}[1]
      \Require taxa $\alpha$, limiar de clipping $c$, batch $m$, passos do crítico $n_{\text{critic}}$
      \Require parâmetros iniciais $w_0$ (crítico), $\thetav_0$ (gerador)
      \While{$\thetav$ não convergiu}
        \For{$t = 1, \ldots, n_{\text{critic}}$}
          \State Amostre $\{x^{(i)}\}_{i=1}^{m} \sim P$ (reais)
          \State Amostre $\{z^{(i)}\}_{i=1}^{m} \sim p(z)$ (prior)
          \State $g_w \leftarrow \nabla_w \left[\tfrac{1}{m}\sum_i f_w(x^{(i)}) - \tfrac{1}{m}\sum_i f_w(g_{\thetav}(z^{(i)}))\right]$
          \State $w \leftarrow w + \alpha \cdot \mathrm{RMSProp}(w, g_w)$
          \State $w \leftarrow \mathrm{clip}(w, -c, c)$
        \EndFor
        \State Amostre $\{z^{(i)}\}_{i=1}^{m} \sim p(z)$
        \State $g_{\thetav} \leftarrow -\nabla_{\thetav}\, \tfrac{1}{m}\sum_i f_w(g_{\thetav}(z^{(i)}))$
        \State $\thetav \leftarrow \thetav - \alpha \cdot \mathrm{RMSProp}(\thetav, g_{\thetav})$
      \EndWhile
    \end{algorithmic}
    \vspace{0.3em}
    \tiny Defaults do paper: $\alpha = 5 \times 10^{-5}$, $c = 0.01$,
    $m = 64$, $n_{\text{critic}} = 5$.
    Adaptado de \citeauthoryear{arjovsky2017wasserstein}.
    \note{
      Pontos a destacar:
      \begin{itemize}
        \item $n_{\text{critic}} = 5$: atualizamos o crítico mais vezes
        que o gerador, já que queremos que ele se aproxime do ótimo.
        \item \emph{Weight clipping}: $w \leftarrow \text{clip}(w,-c,c)$
        impõe Lipschitz de forma grosseira, por isso WGAN-GP usa
        penalização de gradiente.
        \item Usa RMSProp por ser mais estável que Adam nesse cenário
        (segundo o paper).
      \end{itemize}
    }
  \end{frame}

  %--- Slide: WGAN-GP ---
  \begin{frame}{\textit{WGAN-GP} -- Gradient Penalty\footfullcite{gulrajani2017improved}}
    \emph{Weight clipping} impõe Lipschitz, mas reduz a capacidade do
    crítico (pesos saturam em $\pm c$). \textbf{WGAN-GP} substitui o
    clipping por um termo de penalidade na loss:
    \vspace{0.3em}
    \[
      L_{\text{GP}} =
      \underbrace{\sum_{j} f[g(z^{(j)}),\phi] - \sum_{i} f[x^{(i)},\phi]}_{\text{loss WGAN original}}
      + \lambda\,\mathbb{E}_{\hat{x}}\!\left[\,\left(\|\nabla_{\hat{x}} f[\hat{x},\phi]\|_{2} - 1\right)^{2}\,\right]
    \]
    \vspace{0.2em}
    \begin{itemize}
      \item $\hat{x} = \alpha\,x + (1-\alpha)\,x^{*}$ com
      $\alpha \sim \mathcal{U}(0,1)$: pontos interpolados entre real e
      gerada.
      \item Penaliza norma do gradiente \emph{diferente de 1}, isto é,
      Lipschitz é atingida na intensidade certa, não no extremo.
      \item Permite voltar a usar \textsc{Adam} e remover BatchNorm do
      crítico. $\lambda = 10$ funciona bem na prática.
      \item Hoje é o default quando se fala em ``WGAN''.
    \end{itemize}
  \end{frame}

  %======================================================================
  % SECTION 8: TRICKS DE TREINAMENTO
  %======================================================================
  \section{Tricks de Treinamento}

  %--- Slide: Progressive growing ---
  \begin{frame}{\textit{Trick 1}: \textit{Progressive growing}\footfullcite{karras2018progressive}}
    \begin{columns}[T]
      \begin{column}{0.48\textwidth}
        \begin{itemize}
          \item Começamos treinando GAN em resolução baixa
          ($4{\times}4$).
          \item À medida que o treino avança, vamos \textbf{adicionando
          camadas} para dobrar a resolução até $1024{\times}1024$.
          \item Transições suaves entre resoluções evitam instabilidade.
          \item Resultado: imagens de altíssima qualidade (faces
          CelebA-HQ).
        \end{itemize}
      \end{column}
      \begin{column}{0.50\textwidth}
        \centering
        \includegraphics[width=\linewidth]{slides/img/05-geradores/progressive-growing.png}
      \end{column}
    \end{columns}
  \end{frame}

  %--- Slide: Minibatch discrimination ---
  \begin{frame}{\textit{Trick 2}: \textit{Minibatch discrimination}\footfullcite{salimans2016improved}}
    \begin{itemize}
      \item Adicione \textbf{estatísticas do \textit{minibatch}} como
      entrada extra do discriminador.
      \item Exemplos: variância intra-batch, distâncias par-a-par entre
      amostras.
      \item \textbf{Efeito}: força o batch gerado a ter variabilidade
      semelhante ao batch real.
      \begin{itemize}
        \item Ataca diretamente o colapso de moda: amostras muito
        parecidas passam a ser detectáveis.
      \end{itemize}
    \end{itemize}
  \end{frame}

  %--- Slide: Truncation ---
  \begin{frame}{\textit{Trick 3}: \textit{Truncation}\footfullcite{brock2019large}}
    \begin{columns}[T]
      \begin{column}{0.42\textwidth}
        \begin{itemize}
          \item Na hora de \emph{amostrar}, só aceite $z$ cujo módulo
          seja menor que um limiar $\tau$.
          \item $\tau$ grande → mais diversidade, menos qualidade.
          \item $\tau$ pequeno → mais qualidade, menos diversidade.
          \item Permite controlar o trade-off qualidade/diversidade
          \emph{sem retreinar} a rede.
        \end{itemize}
      \end{column}
      \begin{column}{0.56\textwidth}
        \centering
        \includegraphics[width=\linewidth]{slides/img/05-geradores/truncation-trick.png}
      \end{column}
    \end{columns}
  \end{frame}

  %======================================================================
  % SECTION 9: RESUMO E TRILEMA
  %======================================================================
  \section{Panorama Geral}

  %--- Slide: Avaliação ---
  \begin{frame}{Como avaliar uma GAN?}
    Não há log-verossimilhança para reportar. Métricas usuais:
    \begin{itemize}
      \item \textbf{Inception Score}\footfullcite{salimans2016improved}
      \[
        \mathrm{IS} = \exp\!\left(\mathbb{E}_{x^{*}}\!\left[\,D_{KL}\!\left(p(y\mid x^{*})\,\|\,p(y)\right)\,\right]\right)
      \]
      Recompensa imagens de classe nítida ($p(y\mid x^{*})$ concentrada)
      \emph{e} diversidade entre amostras ($p(y)$ espalhada).
      Limitação: depende do classificador (Inception em ImageNet).
      \item \textbf{Fréchet Inception Distance (FID)}\footfullcite{heusel2017fid}:
      ajusta gaussianas em \emph{features} do Inception e calcula
      \[
        \mathrm{FID} = \|\mu_{r}-\mu_{g}\|_{2}^{2} +
        \mathrm{tr}\!\left(\Sigma_{r}+\Sigma_{g}-2\sqrt{\Sigma_{r}\Sigma_{g}}\right)
      \]
      Mais robusto que IS, é o \emph{padrão atual}; valores menores são
      melhores.
    \end{itemize}
    \begin{block}{Regra de bolso}
     Ambas as métricas são sensíveis ao tamanho da amostra.
    \end{block}
  \end{frame}

  %--- Slide: TLDR ---
  \begin{frame}{\textit{TL;DR} sobre GANs}
    \begin{columns}[T]
      \begin{column}{0.48\textwidth}
        \begin{itemize}
          \item Existem diversas arquiteturas com algum controle sobre
          as imagens geradas:
          \begin{itemize}
            \item \textit{Conditional GANs}
            \item \textit{StyleGAN}, \textit{BigGAN}, \ldots
          \end{itemize}
          \item \textbf{Pontos fortes}
          \begin{itemize}
            \item Resultados bastante realistas.
            \item Inferência ``rápida'' (1 passagem pelo gerador).
          \end{itemize}
          \item \textbf{Pontos fracos}
          \begin{itemize}
            \item Difíceis de treinar.
            \item Baixa cobertura de possibilidades (\emph{mode
            collapse}).
          \end{itemize}
        \end{itemize}
      \end{column}
      \begin{column}{0.50\textwidth}
        \centering
        \includegraphics[width=\linewidth]{slides/img/05-geradores/gan-pipeline.png}
      \end{column}
    \end{columns}
  \end{frame}

  %--- Slide: Trilema ---
  \begin{frame}{O trilema dos modelos geradores}
    \begin{columns}[T]
      \begin{column}{0.52\textwidth}
        Todo modelo gerador precisa trocar três objetivos:
        \begin{itemize}
          \item \textbf{Qualidade} das amostras.
          \item \textbf{Cobertura} / diversidade de modos.
          \item \textbf{Velocidade} de amostragem.
        \end{itemize}
        \vspace{0.3em}
        \begin{block}{Posicionamento}
          \begin{itemize}
            \item \textit{GANs}: qualidade + velocidade.
            \item \textit{Diffusion Models}: qualidade + cobertura.
            \item \textit{VAEs}/\textit{Normalizing Flows}: velocidade +
            cobertura.
          \end{itemize}
        \end{block}
      \end{column}
      \begin{column}{0.46\textwidth}
        \centering
        \resizebox{\linewidth}{!}{\generativetrilemma[gan]}
      \end{column}
    \end{columns}
    \note{
      Esse diagrama vai aparecer novamente nas aulas de
      \emph{Normalizing Flows} e \emph{Diffusion Models}, servindo como
      mapa geral da unidade.
    }
  \end{frame}

  %======================================================================
  % REFERÊNCIAS E LEITURAS
  %======================================================================
  \section*{Referências e Leituras}

  \begin{frame}{Leituras Recomendadas}
    \leiturasize
    \begin{itemize}
      \item Capítulo 15: \fullcite{prince2023understanding}
      \item \fullcite{goodfellow2014generative}
      \item \fullcite{radford2015dcgan}
      \item \fullcite{arjovsky2017wasserstein}
      \item \url{https://vincentherrmann.github.io/blog/wasserstein/}
    \end{itemize}
  \end{frame}

  \begin{frame}[allowframebreaks]{Referências}
    \sloppy
    \printbibliography[heading=none]
  \end{frame}

\end{document}
